\begin{thm}


Если поле задано с помощью отображения $y(x)$ в тривиaлизации, 


порожденной картой $x^\alpha$,


то операции дифференцирования Ли и частного


дифференцирования $\partial_\alpha$ поля $y(x)$ коммутируют друг с другом.


\end{thm}





{\it Производная Ли сечения $\wt y(x)$ в репере} 


строится следующим образом. При 


преобразовании $x'=x+t\xi$ элемент $(x-t\xi$, $e(x-t\xi)$, $\wt y(x-t\xi))$ 


переходит в элемент


$$


(x,e-tL_\xi e,\ \wt y(x-t\xi))=(x,\ e(1+tK),\ \wt y(x-t\xi)),


$$


где $K=-e^{-1}L_\xi e$. Полю $\wt y(x-t\xi)$ в репере $e(1+tK)$ 


соответствует поле


$$


\wt y'(x)=f(\wt y(x-t\xi),1+tK)=\wt y(x)-t\bigg(\partial_\xi\wt y-\SP\bigg[ 


\frac{\partial f(y,I_n)}{\partial A}K\bigg]\bigg)+\dotsb


$$


в репере $e(x)$. Производная Ли $L_\xi\wt y$, определяемая как


$$


L_\xi\wt y=\lim_{t\to 0}\frac {\wt y(x)-\wt y'(x)}{t},


$$


равна


$$


L_\xi\wt y=\partial_\xi\wt y-


\SP\bigg[ \frac{\partial f(y,I_n)}{\partial A}K\bigg].


$$


Когда репер $e^\alpha_a\partial_\alpha$


превращается в натуральный репер $(\partial x)=(\partial_1,


\dotsc,\partial_n)$, матрица $K$ становится равной


$\frac{\partial\xi}{\partial x}$ и производная Ли в натуральном репере 


представляет производную Ли в координатах


$$


L_\xi y=\partial_\xi y-\SP


\bigg[ \frac{\partial f(y,I_n)}{\partial A}\frac{\partial\xi}{\partial x}


\bigg].


$$





{\it Производная Ли спиноров}. Производная Ли


спиноров в натуральном репере определяется в виде


\begin{gather*}


L_\xi g=\partial_\xi g+\bigg(\frac{\partial \xi}{\partial x}\bigg)^Tg(x)+


g(x)\frac{\partial \xi}{\partial x}, \\


%


L_\xi\psi=\partial_\xi\psi-\frac{1}{4}\gamma_T\bigg(s^{-1}


\frac{\partial s}{\partial g}\bigg(\bigg(\frac{\partial \xi}


{\partial x}\bigg)^Tg(x)+


g(x)\frac{\partial \xi}{\partial x}\bigg)+s^{-1}


\frac{\partial \xi}{\partial x}s\bigg)\gamma\psi.


\end{gather*}


Исключая производную по направлению $\partial_\xi\psi$


с помощью ковариантной производной $\nabla_\xi\psi$,


производную Ли $L_\xi\psi$ можно записать в виде


$$


L_\xi\psi=\nabla_\xi\psi


-\frac{1}{4}\gamma_Ts^{-1}\bigg(\frac{\partial s}{\partial g}L_\xi g+


\nabla\xi s\bigg)\gamma\psi,\ \ \nabla\xi=(\nabla_\alpha\xi^\beta).


$$


Аналогичным образом находится производная Ли спинора в произвольном репере


\begin{gather*}


L_\xi \wt g=(\nabla \xi)^T\wt g+\wt g\nabla \xi, \\


%


L_\xi\psi=\partial_\xi\wt\psi+\frac{1}{4}\gamma_t\bigg(


-s^{-1}\frac{\partial s}{\partial g}e^{-1}L_\xi e+s^{-1}e^{-1}L_\xi es\bigg)


\gamma \wt\psi=\nabla_\xi\wt\psi-\frac{1}{4}\gamma_Ts^{-1}\bigg(


\frac{\partial s}{\partial\wt g} L_\xi\wt g+\nabla\xi s\bigg)\gamma\wt\psi,


\end{gather*}


где $\nabla\xi=(\nabla_b\xi^a)=(e^a_\beta\nabla_\alpha\xi^\beta e^\alpha_b)=


e^{-1}\nabla \xi e$. При получении связи между производной Ли и ковариантной 


производной необходимо учесть, что $K=\nabla\xi-\omega(\xi)$.





{\it Преобразование производной Ли спинора при


преобразовании координат}.





\begin{thm}


Если в сечении $s: g\to s(g)\in GL(n)$ преобразование


$s(g)$ есть $\eta$-симметрическое


преобразование, то производная Ли спинора есть спинор при


преобразованиях координат.


\end{thm}





\begin{pf}


Так как $\frac{\partial }{\partial \psi}(\Lambda(L)\psi)L_\xi\psi


=\Lambda(L)L_\xi\psi$, то для доказательства теоремы


достаточно доказать, что


$\frac{\partial }{\partial g}\Lambda(L)=0$ или что


$\delta L=\frac{\partial L}{\partial g}\delta g=0$,


где $\delta g$ --- некоторая произвольная вариация


метрики. Так как $L=s^{-1}(g'(x'))\frac{\partial x'}


{\partial x}s(g(x))$, то уравнение $\delta L=0$


эквивалентно уравнению


$$


s^{-1}(g') \frac{\partial s(g')}{\partial g'}\delta g'=


Ls^{-1}(g) \frac{\partial s(g)}{\partial g}\delta g L^{-1}.


$$


Если $s(g)$ --- $\eta$-симметрическое


преобразование, то $s^2(g)=g^{-1}\eta$, откуда


$$


s^{-1}(g) \frac{\partial s(g)}{\partial g}\delta g+


\frac{\partial s(g)}{\partial g}\delta gs^{-1}(g)=


-s\eta^{-1}\delta g s.


$$


Так как относительно преобразования $s(g)$ предполагается, что у него 


собственные значения положительны, а комлексные собственные значения имеют 


положительные вещественные части, то снова можем утверждать, что


последнее уравнение имеет решение


$\frac{\partial s(g)}{\partial g}\delta g$ и оно единственно.


Поэтому для доказательства соотношения


$\delta L=0$ достаточно доказать, что


$$


s'\eta^{-1}\delta g's'=Ls\eta^{-1}\delta gsL^{-1}.


$$


Но это соотношение есть тождество в силу соотношений


$\delta g'=\frac{\partial x}{\partial x'}^T\delta g


\frac{\partial x}{\partial x'}$,


$\frac{\partial x'}{\partial x}=s'Ls^{-1}$.


\end{pf}





Приведем другое доказательство теоремы на основе иного подхода и на


языке реперов.





{\it Закон преобразования производной Ли спинора при 


преобразовании реперов.}


Пользуясь понятием ковариантной производной спинора, производную


Ли спинора можно записать в виде


$$


L_\xi\psi=\bigtriangledown_\xi\psi+\frac{1}{4}\gamma_TS^{-1}(\delta s


-\bigtriangledown\xi s)\gamma\psi,


$$


где


$$


s^{-1}\delta s+\delta ss^{-1}=s\eta^{-1}L_\xi gs,\


\bigtriangledown\xi=(\bigtriangledown_a\xi^b).


$$


Для производной Ли метрического тензора $g$ можно также


использовать формулу


$$


L_\xi g_{ab}=\bigtriangledown_a\xi_b+\bigtriangledown_b\xi_a.


$$


Поскольку $\bigtriangledown_\xi\psi$ --- спинор, то для


доказательства того, что $L_\xi\psi$ --- спинор, необходимо доказать, что


$$


\gamma_T(s'{})^{-1}(\delta s'-\bigtriangledown\xi's')\gamma\psi'=


\Lambda\gamma_Ts^{-1}(\delta s-\bigtriangledown\xi s)\gamma\psi,


$$


если $\psi'=\Lambda\psi$, или при условии, что


$s^{-1}(\delta s-\bigtriangledown\xi s)$ ---


$\eta$-кососимметрическое преобразование, доказать равенство


\begin{equation}


(s'{})^{-1}\delta s'=Ls^{-1}\delta sL^{-1}.


\end{equation}


Воспользовавшись соотношениями


$$


L\gamma=\Lambda^{-1}\gamma\Lambda,\ \ \gamma_TL^{-1}=


\Lambda^{-1}\gamma_T\Lambda,\ \ \bigtriangledown\xi'=A\bigtriangledown\xi


A^{-1},


$$


для доказательства (1) нужно


показать, что определяемое из (1) выражение для $\delta S'$


удовлетворяет уравнению


$$


(s'{})^{-1}\delta s'+\delta s'(s'{})^{-1}=s'\eta^{-1}L_\xi g's'.


$$


Из (1) имеем


$$


\delta s'(s'{})^{-1}=((s'{})^{-1}\delta s')^*=L\delta ss^{-1}L^{-1},


$$


поэтому проверка выполнения предыдущего уравнения сводится к


проверке выполнения соотношения


$$


s'\eta^{-1}L_\xi g's'=Ls\eta^{-1}L_\xi g sL^{-1}.


$$


Это соотношение есть следствие соотношений


$$


L_\xi g=A^TL\xi g'A,\ \ L=(s'{})^{-1}As.


$$


Таким образом, (1) тождественно выполняется, а это означает,


что производная Ли спинора есть спинор относительно


произвольных преобразований реперов.





Ввиду важности теоремы о коммутативности операций дифференцирования Ли 


и частного дифференцирования в общей теории относительности для


построения метрического тензора энергии-импульса проведем


доказательство этой теоремы в частном случае спинорных полей,


рассматриваемых в координатах.





{\it Производная Ли частной производной спинора}. Закон преобразования 


спиноров $\psi(x)$ при преобразовании координат


$x'=x'(x)$ имеет вид $\psi'(x')=\Lambda(L)\psi$,


$L=s^{-1}(g'(x'))\frac{\partial x'}{\partial x}


s(g(x))$, поэтому закон преобразования для частной


производной спинора записывается в виде


$$


\partial_{\alpha'}\psi'(x')=\frac{\partial \Lambda}


{\partial L} \partial _{\alpha'}L\psi(x)+


\Lambda(L)\partial_\beta\psi(x)


\frac{\partial x^\beta}{\partial x'{}^\alpha}.


$$


Если $x'=x+t\xi$, то


$\wt g(x)\equiv g'(x-t\xi)=g(x)-tL_\xi g$,


$\widetilde {\partial_\alpha g(x)}=\partial_\alpha-t


\partial_\alpha L_\xi g$. Тогда $\wt L$, опрeделяемое как


$$


\wt L=s^{-1}(\wt g(x))\frac{\partial (x+t\xi)}{\partial x}s(g(x-t\xi)),


$$


равно


$$


\wt L=1+t L_\xi L,\ \ L_\xi L=s^{-1}\bigg[


\frac{\partial s}{\partial g}(L_\xi g-\partial_\xi g)+


\frac{\partial \xi}{\partial x}s\bigg].


$$


Производную Ли спинора можно записать в виде


$L_\xi\psi=\partial_\xi\psi-\frac{1}{4}\gamma_TL_\xi L\gamma \psi$. Для


$\widetilde{ \partial_\alpha \psi(x)}$ имеем


$$


\widetilde {\partial_\alpha \psi(x)}=\frac{1}{4}\gamma_T


\widetilde {\partial_\alpha L} \gamma\psi(x-t\xi)+


\bigg\{1+\frac{t}{4}\bigg(L_\xi L\bigg)\gamma\bigg\}


\partial_\beta\psi(x-t\xi)


(\delta^\beta_\alpha-t\partial_\alpha\xi^\beta),


$$


где


\begin{multline*}


\widetilde {\partial_\alpha L}=


\frac{\partial s^{-1}(\wt g(x))}{\partial g}


\widetilde {\partial_\beta g}\bigg(1+t\frac{\partial\xi}{\partial x}\bigg)


s(g(x-t\xi)) (\delta^\beta_\alpha-t\partial_\alpha\xi^\beta)+ts^{-1}


\partial_\alpha \frac{\partial\xi}{\partial x}s+ \\


%


+s^{-1}(\wt g(x))\bigg(1+t\frac{\partial\xi}{\partial x}\bigg)\partial_\beta


s(x-t\xi)(\delta^\beta_\alpha-t\partial_\alpha\xi^\beta).


\end{multline*}


Так как


\begin{multline*}


\frac{\partial s^{-1}(\wt g(x))}{\partial g}


\widetilde {\partial_\alpha g(x)}=-s^{-1}(\wt g(x))


\frac{\partial s(\wt g(x))}{\partial g}


\widetilde {\partial_\alpha g(x)}s^{-1}(\wt g(x))= \\


%


=-\bigg(s^{-1}+ts^{-1}\frac{\partial s}{\partial g}L_\xi g s^{-1}\bigg)


\bigg(\frac{\partial sS}{\partial g}-


t\frac{\partial^2 s}{\partial g^2}L_\xi g\bigg)


(\partial_\alpha g-t\partial_\alpha L_\xi g)


\bigg(s^{-1}+ts^{-1}\frac{\partial s}{\partial g}L_\xi g s^{-1}\bigg)= \\


%


= -s^{-1}\frac{\partial s}{\partial g}\partial_\alpha g s^{-1}+tP_1,


\end{multline*}


где


$$


P_1=s^{-1}\bigg\{-\frac{\partial s}{\partial g}L_\xi g


s^{-1}\frac{\partial s}{\partial g}\partial_\alpha g+


\frac{\partial^2 s}{\partial g^2}L_\xi g \partial_\alpha g+


\frac{\partial s}{\partial g}\partial_\alpha L_\xi g-


\frac{\partial s}{\partial g}\partial_\alpha gs^{-1}


\frac{\partial s}{\partial g}L_\xi g\bigg\}s^{-1},


$$


то для первого слагаемого в выражении для


$\widetilde {\partial_\alpha L}$ имеем значения


$$


-s^{-1}\frac{\partial s}{\partial g}\partial_\alpha g+t


\bigg\{P_1s+s^{-1}\frac{\partial s}{\partial g}\partial_\alpha gs^{-1}


\bigg(\frac{\partial s}{\partial g}\partial_\xi g-


\frac{\partial \xi}{\partial x}s\bigg)+


s^{-1}\frac{\partial s}{\partial g}\partial_\beta g


\frac{\partial \xi^\beta}{\partial x^\alpha}\bigg\}.


$$


Третье слагаемое в выражении для


$\widetilde {\partial_\alpha L}$ равно


\begin{multline*}


\bigg(s^{-1}+ts^{-1}\frac{\partial s}{\partial g}L_\xi g s^{-1}\bigg)


\bigg(1+t\frac{\partial \xi}{\partial x}\bigg)


\bigg(\frac{\partial s}{\partial g} \partial_\beta g-\bigg[


\frac{\partial^2 s}{\partial g^2}\partial_\beta g \partial_\xi g


+\frac{\partial s}{\partial g}\partial_\beta \partial_\xi g


\bigg]\bigg)(\delta^\beta_\alpha-t\partial_\alpha\xi^\beta)= \\


%


=s^{-1}\frac{\partial s}{\partial g}\partial_\alpha g+


ts^{-1}\bigg\{\frac{\partial s}{\partial g}L_\xi g


s^{-1}\frac{\partial s}{\partial g}\partial_\alpha g+


\frac{\partial s}{\partial g}\partial_\alpha g- 


\frac{\partial^2 s}{\partial g^2}\partial_\alpha g\partial_\xi g-


\frac{\partial s}{\partial g}\partial_\alpha \partial_\xi g-


\frac{\partial s}{\partial g}\partial_\beta g


\partial_\alpha\xi^\beta\bigg\}.


\end{multline*}


Значит


$$


\widetilde {\partial_\alpha L}=ts^{-1}\bigg\{


\frac{\partial^2 s}{\partial g^2}(L_\xi g-\partial_\xi g)


\partial_\alpha+


\frac{\partial s}{\partial g}\partial_\alpha (L_\xi g-\partial_\xi g)-


\frac{\partial s}{\partial g}\partial_\alpha gL_\xi g+


\partial_\alpha \frac{\partial \xi}{\partial x} s+


\frac{\partial \xi}{\partial x}


\frac{\partial s}{\partial g}\partial_\alpha g\bigg\}=t


\partial_\alpha L_\xi L,


$$


поэтому


$$


\widetilde {\partial_\alpha \psi}= \frac{t}{4}\gamma_T


\partial_\alpha L_\xi L \gamma\psi(x)+


\partial_\alpha \psi(x)+t\bigg(\frac{1}{4} \gamma_TL_\xi L\gamma


\partial_\alpha \psi-\partial_\xi\partial_\alpha\psi-


\partial_\beta\psi\partial_\alpha\xi^\beta\bigg),


$$


откуда


$$


L_\xi(\partial_\alpha\psi)=\partial_\alpha\partial_\xi\psi-


\frac{1}{4}\partial_\alpha(\gamma_TL_\xi L\gamma\psi)=


\partial_\alpha L_\xi\psi.


$$





Доказана





\begin{thm}


Производная Ли частной производной


от спинора равна частной производной от производной Ли спинора.


\end{thm}





\section{Тензор энергии-импульса спинорных полей}





В общей теории относительности функция Лагранжа свободного спинорного поля


в римановом пространстве-времени имеет вид


$$


\cal L=i(\ov \psi\gamma^a\nabla_a\psi-\nabla_a\ov \psi\gamma^a\psi)/2-m


\ov \psi \psi.


$$


Так как при рассмотрении спиноров как в произвольных реперах, так и в 


координатах метрический тензор и гамма-матрицы теряют свой канонический 


вид, то при переходе к лагранжианам произвольной структуры в их аргументы 


нужно включать и метрический тензор, и 


матрицы Дирака. Значит, лагранжиан общего вида имеет вид


$$


\cal L(\psi,\nabla_\alpha\psi,\ov\psi,\nabla_\alpha\ov \psi,\gamma_{g\alpha},


g_{\alpha\beta}).


$$


 


Производная Ли спиноров в координатах обладает свойством:


производная Ли от частной производной есть частная производная


от производной Ли. Из этого свойства следует





\begin{lem}


Тензор $ T^{\alpha\beta}$, 


определяемый с помощью вариационной производной по метрическому


тензору от функции Лагранжа


$$


T^{\alpha\beta}=\frac{2}{\sqrt{-g}}


\frac{\delta( \cal L\sqrt{-g})}{\delta g_{\alpha\beta}},


$$


есть сохраняющийся тензор.


\end{lem}





\begin{pf}


Рассмотрим $\cal L(\psi,\bigtriangledown_\alpha\psi, \ov\psi


\bigtriangledown_\alpha\ov\psi, \gamma_{g\alpha}, g_{\alpha\beta})$ 


как функцию от $g_{\alpha\beta}$, $\partial_\gamma g_{\alpha\beta}$,


$\psi $, $\partial_\alpha\psi$. Tогда при


предположении, что уравнения поля


$\frac{\delta( \cal L\sqrt{-g})}{\delta \psi}=0$


выполнены, имеем


\begin{multline*}


\delta s= \int_\Omega\bigg(


\frac{\partial ( \cal L\sqrt{-g})}{\partial g_{\alpha\beta}}L_\xi


g_{\alpha\beta}+


\frac{\partial ( \cal L\sqrt{-g})}{\partial \partial_\gamma


g_{\alpha\beta}}L_\xi\partial_\gamma g_{\alpha\beta}+


\frac{\partial ( \cal L\sqrt{-g})}{\partial \psi}L_\xi\psi+


\frac{\partial ( \cal L\sqrt{-g})}{\partial \partial_\alpha\psi}


L_\xi\partial_\alpha\psi\bigg)d^4x= \\


%


=\int_\Omega\frac{\delta( \cal L\sqrt{-g})}{\delta g_{\alpha\beta}}


L_\xi g_{\alpha\beta}d^4x=\int_\Omega T^{\alpha\beta}


\bigtriangledown_\beta\xi_\alpha\sqrt{-g}d^4x= 


\int_\Omega\bigtriangledown_\beta T^{\alpha\beta}\xi_\alpha


\sqrt{-g}d^4x=0.


\end{multline*}


Так как область $\Omega$ и векторное поле


$\xi^\alpha $ произвольны, то


$\bigtriangledown_\beta T^{\alpha\beta}=0$.


\end{pf}





Покажем, что тензор $T^{\alpha\beta}$  допускает представление


Л.\,Розенфельда в виде суммы канонического тензора энергии-импульса


и дивергенции тензора углового момента.





Нам понадобятся производные Ли аргументов


$\psi$, $\bigtriangledown_\alpha\psi$, $\ov\psi$,


$\bigtriangledown_\alpha\ov\psi$, $\gamma_{g\alpha}$, $g_{\alpha\beta}$


функции Лагранжа $\cal L(\psi,\bigtriangledown_\alpha\psi, \ov\psi


\bigtriangledown_\alpha\ov\psi, \gamma_{K\alpha}, g_{\alpha\beta})$.


Полагая $P=s^{-1}\frac{\partial s}{\partial g}\big(


(\frac{\partial \xi}{\partial x})^Tg+g\frac{\partial \xi}


{\partial x}\big)+s^{-1}\frac{\partial \xi}{\partial x}s$,


имеем


\begin{gather*}


L_\xi\psi=\partial_ \xi\psi-\frac{1}{4}\gamma_TP\gamma\psi,\quad


L_\xi\ov\psi= \partial_ \xi\ov\psi+\frac{1}{4}\ov\psi\gamma_TP\gamma,\\


%


L_\xi\bigtriangledown_\alpha\psi=\partial_ \xi


\bigtriangledown_\alpha\psi-\frac{1}{4}\gamma_TP\gamma


\bigtriangledown_\alpha\psi+\partial_ \alpha\xi^\beta


\bigtriangledown_\beta\psi,\\


%


L_\xi\bigtriangledown_\alpha\ov\psi=\partial_ \xi\partial


\bigtriangledown_\alpha\ov\psi+\frac{1}{4}\bigtriangledown_\alpha


\ov\psi\gamma_TP\gamma+\partial_ \alpha\xi^\beta


\bigtriangledown_\beta\ov\psi, \\


%


L_\xi\gamma_{g\alpha}= \partial_ \xi\gamma_{g\alpha}+


\frac{1}{2}(s^{-1}{})^\beta_\alpha(P_{\beta\nu}\gamma^\nu-


P^\mu_\beta\gamma_\mu)+\partial_ \alpha\xi^\beta\gamma_{g\beta}.


\end{gather*}


В выкладках Л.\,Розенфельда важную роль играют коэффициенты при


$\partial_ \alpha\xi^\beta$


в выражениях для производной Ли некоторого объекта


$Q$. Эти коэффициенты обозначим $C(Q)^\alpha_\beta$. Тогда


$$


L_\xi Q=\partial_ \xi Q+C(Q)^\alpha_\beta\partial_\alpha\xi^\beta.


$$


Введем обозначения


$$


s^{\mu|\alpha\beta}_\nu=\frac{\partial s^\mu_\nu}


{\partial g_{\alpha\beta}},\quad


P^{\tau|\nu}_{\delta|\mu}=(s^{-1}{})^\tau_\mu s^\nu_\delta+


2(s^{-1}{})^\tau_\sigma s^{\sigma|\alpha\nu}_\delta g_{\alpha\mu}.


$$


Тогда


\begin{gather*}


C(\psi)^\nu_\mu=-\frac{1}{4}\gamma_\tau


P^{\tau|\nu}_{\delta|\mu}\gamma^\delta\psi,\quad


C(\ov\psi)^\nu_\mu=\frac{1}{4}\ov\psi\gamma_\tau


P^{\tau|\nu}_{\delta|\mu}\gamma^\delta, \\


%


C(\bigtriangledown_\alpha\psi)^\nu_\mu=


\delta^\nu_\alpha\bigtriangledown_\mu\psi-\frac{1}{4}\gamma_\tau


P^{\tau|\nu}_{\delta|\mu}\gamma^\delta


\bigtriangledown_\alpha\psi,\\


%


C(\bigtriangledown_\alpha\ov\psi)^\nu_\mu=


\delta^\nu_\alpha\bigtriangledown_\mu\ov\psi+


\frac{1}{4}\bigtriangledown_\alpha\ov\psi\gamma_\tau


P^{\tau|\nu}_{\delta|\mu}\gamma^\delta, \\


%


C(\gamma_{g\alpha})^\nu_\mu =\frac{1}{2}(\eta_4{}_{\delta


\sigma}\gamma^\gamma-\delta^\gamma_\delta\gamma_\sigma)


(s^{-1}{})^\delta_\alpha P^{\sigma|\nu}_{\gamma|\mu}+


\delta^\nu_\alpha\gamma_{g\beta}, \\


%


C(g_{\alpha\beta})^\nu_\mu=\delta^\nu_\beta g_{\alpha\mu}+


\delta^\nu_\alpha g_{\beta\mu}.


\end{gather*}


Вычислим $L_\xi( \cal L\sqrt{-g})$:


\begin{multline*}


L_\xi( \cal L\sqrt{-g})=\frac{\partial( \cal L\sqrt{-g})}


{\partial g_{\alpha\beta}}L_\xi g_{\alpha\beta}+


\frac{\partial( \cal L\sqrt{-g})}{\partial\partial_\gamma


g_{\alpha\beta}}L_\xi\partial_\gamma g_{\alpha\beta}


+\frac{\partial( \cal L\sqrt{-g})}{\partial\psi}L_\xi\psi+ \\


%


\displaybreak[3]


+\frac{\partial


(\cal L\sqrt{-g})}{\partial\partial_\gamma\psi}L_\xi\partial


_\gamma\psi+L_\xi\ov\psi\frac{\partial( \cal L\sqrt{-g})}


{\partial\ov\psi}+L_\xi\partial_\gamma\ov\psi


\frac{\partial( \cal L\sqrt{-g})}{\partial\partial_\gamma\ov\psi}= \\


%


=\frac{\delta( \cal L\sqrt{-g})}{\delta g_{\alpha\beta}}L_\xi g_{\alpha\beta}


+\frac{\delta( \cal L\sqrt{-g})}{\delta\psi}L_\xi\psi+L_\xi\ov\psi


+\frac{\delta( \cal L\sqrt{-g})}{\delta\ov\psi} + \\


%


+\partial_\gamma\bigg(


\frac{\partial( \cal L\sqrt{-g})}{\partial\partial_\gamma


g_{\alpha\beta} }L_\xi g_{\alpha\beta}+\frac{\partial


( \cal L\sqrt{-g})}{\partial\partial_\gamma\psi}L_\xi\psi+


L_\xi\ov\psi\frac{\partial


( \cal L\sqrt{-g})}{\partial\partial_\gamma\ov\psi}\bigg)= \\


%


=\bigg\{\frac{\delta( \cal L\sqrt{-g})}{\delta g_{\alpha\beta}}


\partial_\nu g_{\alpha\beta}+


\frac{\delta( \cal L\sqrt{-g})}{\delta\psi}\partial_\nu\psi +


\partial_\nu\ov\psi\frac{\delta( \cal L\sqrt{-g})}{\delta\psi}- \\


%


-\partial_\mu\bigg[\frac{\delta( \cal L\sqrt{-g})}


{\delta g_{\alpha\beta}}C( g_{\alpha\beta})^\mu_\nu+


\frac{\delta( \cal L\sqrt{-g})}{\delta\psi}C(\psi)^\mu_\nu +


C(\ov\psi)^\mu_\nu\frac{\delta( \cal L\sqrt{-g})}{\delta\psi}


\bigg]\bigg\}\xi^\nu+ \\


%


+\partial_\sigma\bigg\{\bigg[\frac{\delta( \cal L\sqrt{-g})}


{\delta g_{\alpha\beta}}C( g_{\alpha\beta})^\sigma_\mu+


\frac{\delta( \cal L\sqrt{-g})}{\delta\psi}C(\psi)^\sigma_\mu +


C(\ov\psi)^\sigma_\mu\frac{\delta( \cal L\sqrt{-g})}{\delta\ov\psi}+ \\


%


+\frac{\partial( \cal L\sqrt{-g})}


{\partial \partial_\sigma g_{\alpha\beta}}\partial_\mu g_{\alpha\beta}+


\frac{\partial( \cal L\sqrt{-g})}{\partial\partial_\sigma


\psi}\partial_\mu \psi+\partial_\mu\ov\psi


\frac{\partial( \cal L\sqrt{-g})}{\partial\partial_\sigma\ov\psi}


\bigg]\xi^\mu+ \\


%


+\bigg[\frac{\partial( \cal L\sqrt{-g})}


{\partial \partial_\sigma g_{\alpha\beta}}C(


g_{\alpha\beta})^\nu_\mu +


\frac{\partial( \cal L\sqrt{-g})}{\partial\partial_\sigma


\psi}C( \psi)^\nu_\mu+C(\ov\psi)^\nu_\mu


\frac{\partial( \cal L\sqrt{-g})}{\partial\partial_\sigma\ov\psi}


\bigg]\partial_\nu\xi^\mu\bigg\}.


\end{multline*}


Поэтому для $L_\xi( \cal L\sqrt{-g})-\partial_\alpha


(\cal L\sqrt{-g}\xi^\alpha)$ существует представление


$$


L_\xi( \cal L\sqrt{-g})+\partial_\alpha


(\cal L\sqrt{-g}\xi^\alpha)


=Y_\nu\xi^\nu+\partial_\sigma\big(Y^\sigma_\mu\xi^\mu+R


^{\sigma\nu}{}_\mu\partial_\nu\xi^\mu\big),


$$


где


\begin{align*}


Y_\nu&=\frac{\delta( \cal L\sqrt{-g})}{\delta g_{\alpha\beta}}


\partial_\nu g_{\alpha\beta}+


\frac{\delta( \cal L\sqrt{-g})}{\delta\psi}\partial_\nu\psi +


\partial_\nu\ov\psi


\frac{\delta( \cal L\sqrt{-g})}{\delta\psi}- \\


&\quad -\partial_\mu\bigg[\frac{\delta( \cal L\sqrt{-g})}


{\delta g_{\alpha\beta}}C( g_{\alpha\beta})^\mu_\nu+


\frac{\delta( \cal L\sqrt{-g})}{\delta\psi}C(\psi)^\mu_\nu +


C(\ov\psi)^\mu_\nu\frac{\delta( \cal L\sqrt{-g})}{\delta\psi}\bigg], \\


%


Y^\sigma_\mu&= -\cal L\sqrt{-g}\delta^\sigma_\mu+\bigg[


\frac{\delta( \cal L\sqrt{-g})}


{\delta g_{\alpha\beta}}C( g_{\alpha\beta})^\sigma_\mu+


\frac{\delta( \cal L\sqrt{-g})}{\delta\psi}


C(\psi)^\sigma_\mu + \\


%


&\quad +C(\ov\psi)^\sigma_\mu


\frac{\delta( \cal L\sqrt{-g})}{\delta\ov\psi}+


\frac{\partial( \cal L\sqrt{-g})}


{\partial \partial_\sigma g_{\alpha\beta}}\partial_\mu g_{\alpha\beta}+


\frac{\partial( \cal L\sqrt{-g})}{\partial\partial_\sigma


\psi}\partial_\mu \psi+\partial_\mu\ov\psi


\frac{\partial( \cal L\sqrt{-g})}{\partial\partial_\sigma\ov\psi}\bigg], \\


%


R^{\sigma\nu}{}_\mu&= \frac{\partial( \cal L\sqrt{-g})}


{\partial \partial_\sigma g_{\alpha\beta}}C(


g_{\alpha\beta})^\nu_\mu+


\frac{\partial( \cal L\sqrt{-g})}{\partial\partial_\sigma


\psi}C( \psi)^\nu_\mu+C(\ov\psi)^\nu_\mu


\frac{\partial( \cal L\sqrt{-g})}{\partial\partial_\sigma\ov\psi}.


\end{align*}


Равенство нулю интеграла $\int\limits_\Omega \big(


L_\xi( \cal L\sqrt{-g})+\partial_\alpha


( \cal L\sqrt{-g}\xi^\alpha)\big)d^4x$ 


должно иметь место для любой области $\Omega$ и


для любого векторного поля $\xi^\alpha$, поэтому


должны удовлетворяться тождества


$$


Y_\nu=0,\ \ Y^\nu_\mu+\partial_\sigma R^{\sigma\nu}{}_\mu=0,\ \


R^{\sigma\nu}{}_\mu+R^{\nu\sigma}{}_\mu=0.


$$


В предположении, что уравнения поля


$\frac{\delta \cal L}{\delta\psi}=\frac{\delta \cal L}


{\delta\ov\psi}=0 $ удовлетворены, тождество


$ Y_\nu=0 $ эквивалентно тождеству


$$


\frac{\sqrt{-g}}{2}T^{\alpha\beta}\partial_\nu g_{\alpha\beta}-


\partial_\mu(\sqrt{-g}T_\nu^\mu)=0,


$$


которое в свою очередь эквивалентно тождеству


$$


\bigtriangledown_\mu T_\nu^\mu=0.


$$


Cнова пришли к утверждению доказанной выше леммы


о том, что {\it метрический тензор энергии-импульса есть


сохраняющийся тензор}.





Bычислим $R^{\sigma\nu}{}_\mu$. Так как только


$\bigtriangledown_\alpha\psi$ и


$\bigtriangledown_\alpha\ov\psi$ зависят от частных производных


$\partial_\tau g_{\alpha\beta}$ метрического


тензора $g_{\alpha\beta}$ и частных производных


производных $\partial_\alpha\psi$ и


$\partial_\alpha\ov\psi$ спинорных полей $\psi$ и $\ov\psi$, то


$$


R^{\sigma\nu}{}_\mu=


\frac{\partial( \cal L\sqrt{-g})}{\partial\bigtriangledown_\tau\psi}


\frac{\partial\bigtriangledown_\tau\psi}{\partial\partial_\sigma


g_{\alpha\beta}}C( g_{\alpha\beta})^\nu_\mu+


\frac{\partial(\cal L\sqrt{-g})}{\partial\bigtriangledown_\sigma\psi}


C( \psi)^\nu_\mu+


\frac{\partial\bigtriangledown_\tau\ov\psi}{\partial\partial_\sigma


g_{\alpha\beta}}C( g_{\alpha\beta})^\nu_\mu


\frac{\partial( \cal L\sqrt{-g})}{\partial\bigtriangledown_\tau


\ov\psi}+C(\ov \psi)^\nu_\mu


\frac{\partial( \cal L\sqrt{-g})}{\partial\bigtriangledown_


\sigma\ov\psi}.


$$


Так как $\bigtriangledown_\tau\psi=\partial_\tau\psi+


\frac{1}{4}\gamma_TQ_\tau\gamma\psi$, 


$\bigtriangledown_\tau\ov\psi=\partial_\tau\ov\psi-


\frac{1}{4}\ov\psi\gamma_TQ_\tau\gamma$, где


$Q_\tau=s^{-1}\partial_\tau s+s^{-1}\Gamma_\tau s$, 


то для вычисления


$$


\frac{\partial\bigtriangledown_\tau\psi}{\partial\partial_\sigma


g_{\alpha\beta}}


\frac{\partial\bigtriangledown_\tau\ov\psi}{\partial\partial_\sigma


g_{\alpha\beta}}C( g_{\alpha\beta})^\nu_\mu


$$


нужно найти $\frac{\partial Q_\tau}{\partial\partial_\sigma


g_{\alpha\beta}}$. Имеем


$$


\frac{\partial\Gamma_{k\tau s}}{\partial\partial_ig_{jm}}=


\frac{1}{4}\big(\delta^i_\tau(\delta^j_k\delta^m_s +


\delta^m_k\delta^j_s)+


\delta^i_k(\delta^j_\tau\delta^m_s+\delta^m_\tau\delta^j_s)-


\delta^i_s(\delta^m_k\delta^j_\tau+\delta^j_k\delta^m_\tau)\big),


$$


отсюда


\begin{gather*}


\frac{\partial\Gamma_{k\tau}^\rho}{\partial\partial_ig_{jm}}=


\frac{1}{4}\big(\delta^i_\tau(\delta^j_kg^{m\rho} +


\delta^m_kg^{j\rho})+


\delta^i_k(\delta^j_\tau g^{m\rho}+\delta^m_\tau g^{j\rho})-


g^{i\rho}(\delta^m_k\delta^j_\tau+\delta^j_k\delta^m_\tau)\big), \\


%


\frac{\partial\Gamma_{k\tau}^\rho}{\partial\partial_ig_{jm}}


C(g_{jm})^\nu_\mu=


\frac{1}{2}\big(\delta^i_\tau\big[\delta^\nu_k\delta^\rho_\mu +


g_{k\mu}g^{\nu\rho}\big]+


\delta^i_k\big[\delta^\nu_\tau \delta^\rho_\mu+g_{\tau\mu}


g^{\nu\rho}\big]-


g^{i\rho}\big[\delta^\nu_\tau g_{k\mu}+\delta^\nu_k g_{\tau\mu}\big]\big),\\


%


\gamma_{K\rho}\frac{\partial\Gamma_{k\tau}^\rho}


{\partial\partial_ig_{jm}}C(g_{jm})^\nu_\mu\gamma^k_K =


\frac{1}{2}\big\{2\delta^i_\tau\delta^\nu_\mu+


\delta^\nu_\tau\big[\gamma_{K\mu}\gamma^i_K-\gamma^i_K\gamma_{K\mu}\big]-


g^{\tau\mu}\big[\gamma^\nu_K\gamma^i_K-\gamma^i_K\gamma^\nu_K\big]


\big\}.


\end{gather*}


Обозначим $s^{\alpha\beta}=(s^{\mu|\alpha\beta}_\nu)$,


тогда


\begin{gather*}


\gamma_T\frac{\partial(s^{-1}\partial_\tau s)}


{\partial\partial_ig_{jm}}C(g_{jm})^\nu_\mu\gamma =


2\delta^i_\tau\gamma_{gT}s^{\nu m}g_{m\mu}\gamma, \\


%


\frac{\partial\bigtriangledown_\tau\psi}{\partial\partial_\sigma


g_{\alpha\beta}}C( g_{\alpha\beta})^\nu_\mu=


\frac{1}{4}\big(2\delta^\sigma_\tau\gamma_{gT}


s^{\nu m}g_{m\mu}\gamma+\delta^\sigma_\tau\delta^\nu_\mu+


\delta^\nu_\tau\sigma_{g\mu}{}^\sigma+


g_{\tau\mu}\sigma_{g}^{\nu \sigma}\big)\psi,


\end{gather*}


где $\sigma_g{}^{\nu i}=\frac{1}{2}(\gamma_g^\nu\gamma_g^i-


\gamma_g^i\gamma_g^\nu)$. Выражение


$\gamma_\tau P^{\tau|\mu}_{\delta|\mu}\gamma^


\delta$ в $С(\psi)^\nu_\mu$ можно представить в виде


$$


\gamma_\tau P^{\tau|\nu}_{\delta|\mu}\gamma^


\delta=\gamma_{g\mu}\gamma^\nu_g+2\gamma_gs^{\nu m}


g_{m\mu}\gamma,


$$


поэтому


$$


\frac{\partial( \cal L\sqrt{-g})}{\partial\bigtriangledown_\tau\psi}


\frac{\partial\bigtriangledown_\tau\psi}{\partial\partial_\sigma


g_{\alpha\beta}}C( g_{\alpha\beta})^\nu_\mu+


\frac{\partial(\cal L\sqrt{-g})}{\partial\bigtriangledown_\sigma\psi}


C( \psi)^\nu_\mu=


\frac{1}{4}\big(\Phi^\sigma\sigma_K{}^\nu{}_\mu-


\Phi^\nu\sigma_K{}^\sigma{}_\mu+\Phi_\mu\sigma_K{}^{\nu\mu}\big)\psi,


$$


где


$$


\Phi^\sigma=\frac{\partial( \cal L\sqrt{-g})}


{\partial\bigtriangledown_\sigma\psi}.


$$


Аналогичным образом находим


$$


\frac{\partial\bigtriangledown_\tau\ov\psi}{\partial\partial_\sigma


g_{\alpha\beta}}C( g_{\alpha\beta})^\nu_\mu


\frac{\partial(\cal L\sqrt{-g})}


{\partial\bigtriangledown_\tau\ov\psi}+C(\ov \psi)^\nu_\mu


\frac{\partial(\cal L\sqrt{-g})}


{\partial\bigtriangledown_\sigma\ov\psi}=


\frac{1}{4}\ov\psi\big(\sigma^\sigma_{g\mu}\ov\Phi^\nu-


\sigma^\nu_{g\mu}\ov\Phi^\sigma-


\sigma_g{}^{\nu\sigma}\ov\Phi_\mu\big),


$$


где


$$


\ov\Phi^\sigma=\frac{\partial( \cal L\sqrt{-g})}


{\partial\bigtriangledown_\sigma\ov\psi}.


$$


В итоге для $R^{\sigma\nu}{}_\mu $ получим представление


$$


R^{\sigma\nu}{}_\mu=U^{\sigma\nu}{}_\mu+\ov U^{\sigma\nu}{}_\mu,


$$


где


\begin{align*}


U^{\sigma\nu}{}_\mu&=\frac{1}{4}\big(\Phi^\sigma


\sigma^\nu_{g\mu}-


\Phi^\nu\sigma^\sigma_{g\mu}+\Phi_\mu\sigma_g{}^{\nu\mu}\big)\psi, \\


%


\ov U^{\sigma\nu}{}_\mu &=


\frac{1}{4}\ov\psi\big(\sigma^\sigma_{g\mu}\ov\Phi^\nu-


\sigma^\nu_{g\mu}\ov\Phi^\sigma-


\sigma_g{}^{\nu\sigma}\ov\Phi_\mu\big).


\end{align*}


Осталось найти представление для


$Y^\sigma_\mu$. Для этого уже надо вычислить следующее:


$$


\frac{\partial\bigtriangledown_\tau\psi}


{\partial\partial_\sigma g_{\alpha\beta}}\partial_\mu


g_{\alpha\beta}=2\frac{\partial\bigtriangledown_\tau\psi}


{\partial\partial_\sigma g_{\alpha\beta}}\Gamma^\nu_{\mu\alpha}


g_{\nu\beta}=\frac{1}{2}\gamma_T\frac{\partial Q_\tau}


{\partial\partial_\sigma g_{\alpha\beta}}\Gamma^\nu_{\mu\alpha}


g_{\nu\beta}\gamma\psi=


\frac{1}{2}\gamma_T\bigg(s^{-1}\frac{\partial\partial _\tau s}


{\partial\partial_\sigma g_{\alpha\beta}}+


s^{-1}\frac{\partial


\Gamma_\tau}{\partial\partial_\sigma g_{\alpha\beta}}


s\bigg)\Gamma^\nu_{\mu\alpha} g_{\nu\beta}\gamma\psi.


$$


Обратимся к вычислению


\begin{gather*}


\gamma_Ts^{-1}\frac{\partial


\Gamma_\tau}{\partial\partial_\sigma g_{\alpha\beta}}


s\Gamma^\nu_{\mu\alpha} g_{\nu\beta}\gamma\psi=


\gamma_{g\rho}\frac{\partial


\Gamma^\rho_{k\tau}}{\partial\partial_\sigma g_{\alpha\beta}}


s\Gamma^\nu_{\mu\alpha} g_{\nu\beta}\gamma\psi= \\


%


=\frac{1}{4}\gamma_{g\rho}\big\{\delta^\sigma_\tau\big[


\Gamma^\rho_{\mu k}+\Gamma^\nu_{\mu j}g^{j\rho}g_{\nu k}\big]+


\delta^\sigma_k\big[\Gamma^\rho_{\mu \tau}+\Gamma^\nu_{\mu j}


g_{\nu\tau}g^{j\rho}\big]-


g^{\sigma\rho}\big[\Gamma^\nu_{\mu \tau}g_{\nu k}+


\Gamma^\nu_{\mu k}g_{\nu\tau}\big]\big\}\gamma^k_g\psi.


\end{gather*}


Также находим


$$


\gamma_Ts^{-1}\frac{\partial\partial _\tau s}


{\partial\partial_\sigma g_{\alpha\beta}}


\Gamma^\nu_{\mu\alpha} g_{\nu\beta}\gamma\psi=


\gamma_gs^{\alpha\beta}\delta^\sigma_\tau\Gamma^\nu_{\mu\alpha}


g_{\nu\beta}\gamma\psi.


$$


Далее имеем


\begin{gather*}


\frac{\partial( \cal L\sqrt{-g})}{\partial\bigtriangledown_\tau\psi}


\frac{\partial\bigtriangledown_\tau\psi}{\partial\partial_\sigma


g_{\alpha\beta}}\partial_\mu g_{\alpha\beta}+


\frac{\partial(\cal L\sqrt{-g})}{\partial\bigtriangledown_\sigma\psi}


\partial_\mu\psi= \\


%


=\Phi^\tau\bigg(\frac{1}{8}


\gamma_{g\rho}\big\{\delta^\sigma_\tau\big[


\Gamma^\rho_{\mu k}+\Gamma^\nu_{\mu j}g^{j\rho}g_{\nu k}\big]+


\delta^\sigma_k\big[\Gamma^\rho_{\mu \tau}+\Gamma^\nu_{\mu j}


g_{\nu\tau}g^{j\rho}\big] -\\


%


-g^{\sigma\rho}\big[\Gamma^\nu_{\mu \tau}g_{\nu k}+


\Gamma^\nu_{\mu k}g_{\nu\tau}\big]\big\}\gamma^k_g\psi+


\frac{1}{2}\gamma_{gT}s^{\alpha\beta}\delta^\sigma_\tau


\Gamma^\nu_{\mu\alpha}g_{\nu\beta}\gamma\psi\bigg)+ \\


%


+\Phi^\sigma\bigg(\bigtriangledown_\mu\psi-\frac{1}{4}


\big[2\gamma_{gT}s^{\alpha\beta}\delta^\sigma_\tau


\Gamma^\nu_{\mu\alpha}g_{\nu\beta}\gamma\psi+


\gamma_{gT}\Gamma_\mu\gamma_g\big]\psi\bigg)=


\Phi^\sigma\bigtriangledown_\mu\psi+Z^\sigma_\mu,


\end{gather*}


где


\begin{multline*}


Z^\sigma_\mu=\frac{1}{8}\big\{\Phi^\sigma\big(


\gamma_g^j\Gamma^\nu_{\mu j}\gamma_{g\nu}-\gamma_{KT}


\Gamma_\mu\gamma_g\big)+ \\


%


+\Phi^\tau\gamma_{g\rho}\gamma_{gT}\big(\Gamma^\rho_{\mu\tau}


+\Gamma^\nu_{\mu j}g_{\nu\tau}g^{j\rho}\big)\gamma_g^\sigma-


\Phi^\tau\gamma_g^\sigma\big(\Gamma^\nu_{\mu\tau}g_{\nu k}+


\Gamma^\nu_{\mu j}g_{\nu\tau}\big)\gamma_g^k\big\}\psi.


\end{multline*}


Пусть $R^{\sigma\nu}{}_\mu/\sqrt{-g}=P^{\sigma\nu}{}_\mu$,


где $P^{\sigma\nu}{}_\mu$ является тензором, тогда


$$


\partial_\sigma R^{\sigma\nu}{}_\mu=\sqrt{-g}\bigtriangledown_


\sigma P^{\sigma\nu}{}_\mu+\Gamma^\tau_{\sigma\mu}


R^{\sigma\nu}{}_\tau.


$$


Имеет место


$$


Z^\nu_\mu+\Gamma^\tau_{\sigma\mu}U^{\sigma\nu}{}_\tau=0.


$$


Аналогичными вычислениями находим


$$


\frac{\partial\bigtriangledown_\tau\ov\psi}{\partial\partial_\sigma


g_{\alpha\beta}}\partial_\mu g_{\alpha\beta}


\frac{\partial( \cal L\sqrt{-g})}{\partial\bigtriangledown_\tau


\ov\psi}+\partial_\mu\ov\psi


\frac{\partial(\cal L\sqrt{-g})}{\partial\bigtriangledown_\sigma


\ov\psi}=\bigtriangledown_\mu\ov\psi\ov\Phi^\sigma-


\Gamma^\tau_{\sigma\mu}\ov U^{\sigma\nu}{}_\tau.


$$


Для $Y^\nu_\mu+\partial_\sigma R^{\sigma\nu}{}_\mu$ получим


$$


Y^\nu_\mu+\partial_\sigma R^{\sigma\nu}{}_\mu=-\cal L\sqrt{-g}


\delta^\nu_\mu+\sqrt{-g}T^\nu_\mu+


\Phi^\nu\bigtriangledown_\mu\psi+\bigtriangledown_\mu


\ov\psi\ov\Phi^\nu+\sqrt{-g}


\bigtriangledown_\sigma P^{\sigma\nu}{}_\mu,


$$


откуда


$$


T^\nu_\mu=\cal L\delta^\nu_\mu-


\big(\Phi^\nu\bigtriangledown_\mu\psi+\bigtriangledown_\mu


\ov\psi\ov\Phi^\nu\big)/\sqrt{-g}-


\bigtriangledown_\sigma P^{\sigma\nu}{}_\mu.


$$





Доказана





\begin{thm}


Пусть функция Лагранжа


$\cal L(\psi,\bigtriangledown_\alpha\psi, \ov\psi,


\bigtriangledown_\alpha\ov\psi, \gamma_{K\alpha},


g_{\alpha\beta})$ есть инвариант


относительно произвольных преобразований координат и удовлетворены


уравнения спинорного поля


$\frac{\delta\cal L}{\delta\psi}=0$. Тогда тензор


$T^{\alpha\beta}=\frac{2}{\sqrt{-g}}\frac{\delta\cal L


\sqrt{-g}}{\delta g_{\alpha\beta}}$


есть сохраняющийся тензор и допускает представление


Л.\,Розенфельда в виде суммы канонического тензора энергии-импульса 


и дивергенции тензора углового момента.


\end{thm}





Автор выражает признательность М.А.\,Малахальцеву за


обсуждение результатов работы и ценные критические замечания.








\begin{thebibliography}{99}





\bibitem{1}


Cartan E. {\em Les groups projectifs qui ne laissent invariants aucune 


multiplicite plane} // Bull. Soc. Math. de France, 1913. -- V.\,41.


-- P.\,53--96.


\bibitem{2}


Van der Waerden B.L. {\em Spinoranalyse} // Nachr. Acad. Wiss. G\"ottingen, 


Math. -- Physik. Kl. -- 1929. -- P.\,100--109.


\bibitem{3}


Fock V.A., Ivanenko D.D. {\em Geometrie quantique lineare et deplacement 


parallele} // Compt. Rend. Acad. Sci. Paris, 1929. -- V.\,188.


-- P.\.1470--1472.


\bibitem{4}


Rosenfeld L. {\em Sur le tenseur d'impulsion-energy} // Mem. Acad. Roy. 


Belgique, 1940. -- V.\,18. -- \No\,6. -- P.\,1--30.


\bibitem{5}


Lichnerowicz A. {\em Spineurs harmoniques} // Compt. Rend. Acad. Sci. Paris,


1963. -- V.\,253. -- \No\,1. -- P.\,7--9.


\bibitem{6}


Kosmann Y. {\em D\'eriv\'ees de Lie des spineurs} // Compt. Rend.


Acad. Sci. Paris, 1966. -- V.\,262\,A. -- P.\,289--292.


\bibitem{7}


Bourguignon J.-P., Gauduchon P. {\em Op\'erateurs de Dirac et variations 


de m\'etriques} // Commun. 


Math. Phys. -- 1992. -- V.\,144. -- P.\,581--599.


\bibitem{8}


Билялов Р.Ф. {\em Законы сохранения для спинорных полей на римановых


пространственно-временных многообразиях} // Теор. и матем. физика. --


1992. -- T.\,90. -- \No\,3. -- C.\,369--379.


\bibitem{9}


Билялов Р.Ф. {\em Симметрический тензор энергии-импульса


спинорных полей} // Теор. и матем. физика. -- 1996. --


T.\,108. -- \No\,2. -- C.\,306--314.


\bibitem{10}


Кобаяси Ш., Номидзу К. {\em Основы дифференциальной геометрии}. T.\,1.


-- М.: Наука, 1981. -- 344 с. 


\bibitem{11}


Петров А.З. {\em Новые методы в общей теории относительности}.


-- М.: Наука, 1966. -- 495 c.


\bibitem{12}


Ланкастер П. {\em Теория матриц}. Гл. ред. физ.-мат. лит-ры, 1982. -- 270 с.





\end{thebibliography}





\vskip 2cm





\postup{\it Казанский государственный университет}{\it Поступили}


\postup{}{\it первый вариант $27.05.1998$}


\postup{}{\it окончательный вариант $24.09.2002$}





\label{end}


